% Authoring copy of the extension now incorporated literally into B_rigid7.tex.
% Prepared after checkpoint03; all13 paths and auxiliary finite checks certified.
\subsection*{Explicit embedded form of the rigid test kernel}

The following specifies the homomorphism defining the kernel in B.56
on a geometric free basis, and gives an exact membership procedure
for finite words.  It does not identify the completed kernel with
the closure of its discrete counterpart.

\paragraph{A certified geometric basis.}
Retain the base point $u_*=(371+237i)/100$, radius
$\epsilon=3/250$, and square-meridian paths of B.56.
Put
\[
 A=36603/100000,\qquad B=136603/100000.
\]
The thirteen square centers, in their counterclockwise fan order
as seen from $u_*$, are
\[
 \bigl(-B+Bi,\ i,\ B+Bi,\ -A+Ai,\ -1,\ A+Ai,\ 0,
       -A-Ai,\ -B-Bi,\ A-Ai,\ 1,\ -i,\ B-Bi\bigr).
 \tag{B.56.5}
\]
Let $e_j$ follow the straight stem to the northeast corner of
the $j$th square, traverse its boundary counterclockwise, and
return along the same stem.  These are a free basis of
$\pi_1(S_T,u_*)$.

Here is the exact geometric verification, including the
otherwise essential basis condition.  The thirteen closed
squares are disjoint.  Rational segment--rectangle intersection
tests show that each stem meets its own square only at its
northeast corner and misses every other square.  The thirteen
stems have distinct slopes and a common initial point, so their
interiors are disjoint.  Every square contains precisely its
named exceptional parameter point.  For the irrational points,
this follows from
\[
 \frac{17320508075688772}{10^{16}}<\sqrt3<
 \frac{17320508075688773}{10^{16}},
\]
checked by squaring, and the definitions of $a,b$ in B.56.
All these tests are over $\mathbb Q$.  The northeast ray from
$u_*$ to infinity misses the whole construction.  Cutting along
this ray and the embedded petal tree gives the usual geometric
meridian basis of the thirteen-puncture plane.  Its positively
oriented meridian at infinity is
$e_{14}=(e_1\cdots e_{13})^{-1}$.
In particular the previous certified word is literally
$h=e_6e_{11}e_6^{-1}e_5^{-1}$ in this basis.

\paragraph{All thirteen monodromies.}
Retain the set $S_0$ of B.56 and send labels $2,4,6$ to
$\infty,0,1$.  The six finite labeled points define a map
$S_T\to\operatorname{Conf}_6(\mathbb C)$ and therefore a
specified homomorphism
\[
 \theta:F_{13}\longrightarrow PB_6,\qquad e_j\longmapsto\beta_j.
 \tag{B.56.6}
\]
All thirteen words $\beta_j$ are recorded in
\texttt{all13\_braid\_words.json} and
\texttt{all13\_braid\_words.txt}; the sign convention is that
of B.56.  Their reduced lengths and exponent sums are
\[
\begin{array}{c|rrrrrrrrrrrrr}
j&1&2&3&4&5&6&7&8&9&10&11&12&13\\\hline
|\beta_j|&60&58&26&54&58&36&58&82&88&44&14&72&30\\
\operatorname{exp}(\beta_j)&26&-38&26&-4&2&-4&-8&26&-4&26&2&-38&-4
\end{array}
\]
The program \texttt{certify\_all13\_monodromy.py} checks every
crossing on all $78$ rational path segments by the same exact
root-isolation procedure as B.56.  It excludes collisions,
poles and exceptional parameters on every segment, checks all
endpoint orders, and verifies the identity permutation of
each word.  The three previous words for indices $5,6,11$
agree literally with their original certificate.
The complete polynomial and isolating-interval data are in
\texttt{all13\_exact\_certificate.json}; actual stdout is
\texttt{certify\_all13\_monodromy.out}.

\paragraph{The exact preimage and finite-word membership.}
Let $z=\Delta_6^2=(\sigma_1\cdots\sigma_5)^6$.
Capping the disk at its boundary gives
$PB_6/\langle z\rangle=P\Gamma_{0,7}$, with the cap point
the named puncture $2$; see the capping comparison and center
calculation in \citet[Sections 2.2 and 4.3]{DE_GonzalezMeneses2011}.
There is also a useful splitting check.  Forgetting all but the
fixed strands with labels $4,6$ gives
$\ell_{4,6}:PB_6\to PB_2\cong\mathbb Z$.
Those two normalized points are constantly $0,1$, so
$\ell_{4,6}\theta=0$, whereas $\ell_{4,6}(z)=1$.
Consequently the discrete intersection requested in B.4(a) is
the explicitly specified subgroup
\[
                 N_{\mathrm{rig}}=\theta^{-1}(\langle z\rangle)=\ker\theta.
 \tag{B.56.7}
\]
This gives a terminating exact membership test for every finite
word in the thirteen source generators.  Substitute the thirteen
specified braid words.  In general, centrality is tested by
rejecting exponent sums not divisible by $30$ and otherwise
putting $m$ equal to that sum divided by $30$.
With $v=x_1\cdots x_6$, compare the six freely reduced Artin
images to
\[
                       x_j\longmapsto v^m x_jv^{-m}.
\]
Accept precisely if all six agree.  The Artin action is faithful
and the displayed automorphism is the action of $z^m$, so this
criterion is both necessary and sufficient.  For our normalized
source words, $\ell_{4,6}=0$ forces $m=0$, so the implementation
rejects every nonzero exponent sum and otherwise compares
the Artin images to the six generators themselves.
No efficiency bound
is claimed for this free-word expansion.
\texttt{membership\_rigid7.py} implements it and checks the
full-twist convention, the word $h$, and all fourteen peripheral
meridians.  Its actual output is retained.

Completing (B.56.6), let
$Z_{\mathrm{tw}}=\overline{\langle z\rangle}\leq\widehat{PB_6}$.
This central closure is isomorphic to $\widehat{\mathbb Z}$:
the homomorphism $\ell_{4,6}$ restricts to an isomorphism
$\langle z\rangle\to\mathbb Z$, so it induces the full
profinite topology on this cyclic subgroup.  The universal property of completion
gives
\[
 \widehat{PB_6}/Z_{\mathrm{tw}}\cong\widehat{P\Gamma_{0,7}}.
\]
Indeed finite quotients of the right-hand discrete group are
exactly finite quotients of $PB_6$ killing $z$.
Thus the completed test kernel has the exact embedded description
\[
       N_{\mathrm{rig}}^{\mathrm{test}}
                    =\widehat\theta^{-1}(Z_{\mathrm{tw}})=\ker\widehat\theta.
 \tag{B.56.8}
\]
For the last equality, $\widehat\ell_{4,6}$ is an isomorphism
on $Z_{\mathrm{tw}}$ and vanishes on the image of
$\widehat\theta$.  The independent exact program
\texttt{check\_all13\_pair\_linking.py} checks the initial
label order and all fifteen pair linkings of all thirteen words.
It confirms that the $(4,6)$ linking is zero in every case and
that the full twist has linking one.
The thirteen finite words specify this continuous map completely.
The argument uses the central closure, and does not need a
separate assertion about the full center of $\widehat{PB_6}$.
It supplies no new induced-topology comparison for
$\theta(F_{13})$ or its spherical image.

\paragraph{The rigid and four-point kernels are incomparable.}
There is a short exact test detecting $h$ among the four-point
maps.  For $q=\{1,2,3,5\}$, the cross-ratio of B.18 is
\[
                   \psi_q(u)=\left(\frac{u^2-1}{u^2+1}\right)^2.
 \tag{B.56.9}
\]
It is an unramified degree-four covering
\[
 \mathbb P^1\setminus\{0,\infty,\pm1,\pm i\}
                   \longrightarrow\mathbb P^1\setminus\{0,1,\infty\}.
\]
Its zeros and poles are respectively $\pm1,\pm i$, each of
multiplicity two; the preimages of $1$ are $0,\infty$, also
of multiplicity two.  These account for all ramification.
Filling the other eight source punctures kills $e_6$ and
carries $h$ to the difference of the two distinct meridians
at $1,-1$, nonzero already in homology.  The covering induces
an injection on fundamental groups.  Hence
$\psi_{q*}(h)\ne1$.  Residual finiteness of the target
free group shows that this same finite word survives the
completed four-point test.

In the other direction, form the balanced commutator $w$ of
B.19 using the present fourteen meridians.  Its reduced source
word has length $384$.  It dies in every four-point test.
The auxiliary exact program \texttt{check\_all13\_burau.py}
sends $B_6$ to $\operatorname{GL}_6(\mathbb F_{101})$ by the
unreduced Burau matrices at $t=2$: the block at positions
$j,j+1$ is
$\left(\begin{smallmatrix}-1&2\\1&0\end{smallmatrix}\right)$,
with identity elsewhere.  The program checks all defining
braid relations and all inverse matrices.  It obtains
\[
 \rho(\theta(w))=
 \begin{pmatrix}
 77&19&64&20&74&50\\
 47&60&88&59&42&8\\
 55&6&5&12&79&46\\
 8&98&44&86&18&50\\
 50&0&26&19&92&16\\
 5&62&69&96&93&80
 \end{pmatrix}\ne I.
\]
The exponent sum of $\theta(w)$ is zero, so if it were a
central full-twist power it would be the identity.  This
matrix proves $w\notin N_{\mathrm{rig}}$.  In fact the matrix
fails to commute with every displayed braid generator, so
it also survives the finite braid image modulo the central
image of $z$.  This directly witnesses nontriviality in a
finite quotient of the seven-puncture target.
Thus neither the rigid nor the four-point kernel contains
the other, discretely or after completion.  Their intersection
contains the explicit $776$-letter source word $[h,w]$.
The matrix representation is used only to prove a nonidentity;
its faithfulness is not assumed.
